\section{CHƯƠNG 1. MÔ TẢ BÀI TOÁN VÀ CƠ SỞ LÝ THUYẾT}
\vspace{0.5cm}

\subsection{1. Mô tả chi tiết bài toán}
\subsubsection{1.1. Bài toán phát hiện chữ bị làm giả trong tài liệu}
\indent\indent Nghiên cứu tập trung giải quyết bài toán phát hiện và phân loại các vùng văn bản bị làm giả (Tampered Text Detection) trên các ảnh tài liệu số hóa. Khác với bài toán OCR thông thường chỉ đọc nội dung văn bản, bài toán pháp y tài liệu đòi hỏi mô hình phải phân tích sâu các đặc trưng vi mô ở cấp độ pixel và miền tần số để xác định tính xác thực của từng ký tự hoặc khối văn bản.\vspace{0.3cm}

Các hình thức làm giả văn bản phổ biến trong tài liệu bao gồm:\vspace{-0.1cm}
\begin{itemize}[label={--}, itemsep=1pt]
    \item \textbf{Sao chép - Dán (Copy-Move):} Kẻ giả mạo sao chép một cụm ký tự hoặc con số từ một vị trí trong cùng tài liệu (ví dụ: con số 0 trong ngày tháng) và dán đè vào vị trí khác (như số tiền giao dịch). Dấu hiệu đặc trưng là sự trùng khớp tuyệt đối về độ nhiễu và hoa văn font chữ nhưng lại có viền tiếp giáp bất thường.
    \item \textbf{Ghép nối văn bản (Text Splicing):} Lấy chữ số hoặc từ ngữ từ một tài liệu nguồn khác ghép vào tài liệu đích. Trường hợp này thường để lại các dấu vết sai lệch rõ rệt về độ phân giải, độ sắc nét viền, góc nghiêng hoặc bảng màu nền.
    \item \textbf{Chỉnh sửa hoặc thay thế phông chữ (Digital Font Replacement):} Sử dụng công cụ đồ họa máy tính để gõ lại nội dung chữ mới đè lên nền đã xóa. Dấu hiệu là sự thiếu tự nhiên về độ răng cưa (aliasing), sai khác về khoảng cách ký tự (kerning) và độ tương phản quang sai.
    \item \textbf{Xóa văn bản (Text Erasure / Inpainting):} Sử dụng thuật toán lấp đầy nội dung để xóa bỏ chữ ký, con dấu hoặc các điều khoản quan trọng.
\end{itemize}

\subsubsection{1.2. Hệ thống nhãn phân loại pháp y vi mô}
\indent\indent Để mô hình hóa chi tiết bản chất của vùng chữ bị can thiệp, hệ thống định nghĩa 6 nhãn đặc trưng độc lập:
\begin{itemize}[label={--}, itemsep=2pt]
    \item \textbf{Edge artifacts:} Vết cắt dán, viền răng cưa bất thường, dấu vết mờ nhạt do nội suy xung quanh nét chữ.
    \item \textbf{Unnatural texture:} Độ nhiễu nén JPEG, hạt nhiễu cảm biến hoặc độ nhẵn của giấy nền không đồng nhất với tài liệu gốc.
    \item \textbf{Inconsistent alignment:} Ký tự bị lệch dòng, khoảng cách giữa các chữ cái (kerning) hoặc khoảng cách dòng (leading) không chuẩn.
    \item \textbf{Inconsistent font:} Sử dụng font chữ khác biệt về độ dày nét (weight), độ nghiêng, hoặc kiểu dáng chân chữ (serif).
    \item \textbf{Lack of integration:} Vùng chữ thiếu tính hòa nhập vào nền, độ tương phản ánh sáng và đổ bóng không tự nhiên.
    \item \textbf{Text incoherence:} Nội dung văn bản bị vô nghĩa, mâu thuẫn ngữ nghĩa hoặc sai lệch logic với phần còn lại của văn bản.
\end{itemize}

\subsubsection{1.3. Yêu cầu của kiến trúc học đa nhiệm (Multi-task Learning)}
\indent\indent Một hệ thống giám định toàn diện cần đáp ứng đồng thời hai nhiệm vụ:
\begin{enumerate}
    \item \textbf{Phân đoạn mức điểm ảnh (Pixel-level Segmentation):} Trả về mặt nạ nhị phân chính xác kích thước $H \times W$, chỉ rõ từng pixel thuộc vùng chữ bị làm giả.
    \item \textbf{Phân loại thuộc tính giả mạo (Multi-label Classification):} Dự đoán xác suất cho 6 nhãn đặc trưng vi mô, cung cấp căn cứ giải thích vì sao vùng chữ đó bị coi là giả mạo.
\end{enumerate}
Mô hình đề xuất \textbf{LightweightDTD} được thiết kế để giải quyết đồng thời hai nhiệm vụ này trên một mạng chia sẻ đặc trưng duy nhất.

\subsection{2. Cơ sở lý thuyết liên quan}

\subsubsection{2.1. Mạng nơ-ron tích chập EfficientNet-B0}
\indent\indent Kiến trúc \textbf{EfficientNet} do Tan và Le \cite{tan2019efficientnet} đề xuất giải quyết bài toán tối ưu hóa tài nguyên mạng tích chập thông qua phương pháp \textbf{Co giãn phức hợp (Compound Scaling)}. Thay vì mở rộng tùy ý một chiều đơn lẻ, phương pháp này sử dụng một hệ số phức hợp $\phi$ để co giãn đồng thời cả ba chiều theo tỷ lệ cân bằng toán học:
\begin{equation}
\begin{aligned}
\text{Độ sâu: } & d = \alpha^\phi \\
\text{Độ rộng: } & w = \beta^\phi \\
\text{Độ phân giải: } & r = \gamma^\phi \\
\text{thỏa mãn: } & \alpha \cdot \beta^2 \cdot \gamma^2 \approx 2 \quad \text{với } \alpha \ge 1, \beta \ge 1, \gamma \ge 1
\end{aligned}
\end{equation}

Khối xây dựng cốt lõi của EfficientNet-B0 là \textbf{MBConv (Mobile Inverted Bottleneck Convolution)} kết hợp khối chú ý \textbf{Squeeze-and-Excitation (SE)}. Khối MBConv đầu tiên mở rộng số kênh qua lớp tích chập $1 \times 1$, sau đó áp dụng tích chập phân tách theo chiều sâu (Depthwise Convolution $3 \times 3$ hoặc $5 \times 5$), đi qua khối SE để hiệu chỉnh trọng số động cho từng kênh đặc trưng, và cuối cùng chiếu ngược trở lại số kênh ban đầu bằng tích chập $1 \times 1$ kết hợp liên kết tắt (residual connection). Nhờ cấu trúc này, EfficientNet-B0 chỉ có khoảng 5.3 triệu tham số backbone nhưng đạt hiệu quả trích xuất đặc trưng vượt trội.

\subsubsection{2.2. Phép biến đổi Cosin rời rạc theo khối (BlockDCT $8 \times 8$)}
\indent\indent Biến đổi Cosin rời rạc (Discrete Cosine Transform -- DCT) \cite{ahmed1974discrete} là kỹ thuật biến đổi tín hiệu từ miền không gian sang miền tần số, đóng vai trò nền tảng trong chuẩn nén ảnh JPEG. Trong nén JPEG, ảnh được chia thành các khối độc lập kích thước $8 \times 8$ pixel và được lượng tử hóa theo từng tần số.\vspace{0.3cm}

Để trích xuất đặc trưng tần số phục vụ giám định mà không làm gián đoạn luồng tính toán song song trên GPU, mô-đun \textbf{BlockDCT} xây dựng cơ sở biến đổi DCT-II trực chuẩn $8 \times 8$:
\begin{equation}
B_{u, x} = \alpha(u) \cos \left( \frac{(2x + 1)u\pi}{2n} \right), \quad n=8
\end{equation}
trong đó:
\begin{equation}
\alpha(u) = \begin{cases}
\sqrt{\frac{1}{n}} & \text{khi } u = 0 \\
\sqrt{\frac{2}{n}} & \text{khi } u > 0
\end{cases}
\end{equation}

Khi đưa ảnh xám đầu vào kích thước $[B, 1, H, W]$, phép chiếu trực chuẩn theo hai chiều khối $8 \times 8$ sẽ biến đổi ảnh thành tensor hệ số tần số có kích thước $[B, 64, H/8, W/8]$, trong đó mỗi kênh tương ứng với một thành phần tần số cụ thể (từ DC tần số thấp đến các hệ số AC tần số cao). Bất kỳ sự chỉnh sửa ghép nối nào làm xô lệch cấu trúc $8 \times 8$ gốc đều biểu hiện rõ rệt trên các kênh tần số này.

\subsubsection{2.3. Mô-đun chú ý không gian và kênh đồng thời (scSE Module)}
\indent\indent Mô-đun \textbf{scSE (Concurrent Spatial and Channel Squeeze-and-Excitation)} do Roy và cộng sự \cite{roy2018concurrent} đề xuất giúp tăng cường các đặc trưng không gian và ngữ nghĩa thông qua hai luồng chú ý song song:
\begin{enumerate}
    \item \textbf{Chú ý kênh (Channel Squeeze-and-Excitation -- cSE):} Nén chiều không gian thông qua Adaptive Average Pooling và học trọng số kênh qua hai lớp tích chập $1 \times 1$:
    \begin{equation}
    \mathbf{z}_{cSE} = \sigma \left( \mathbf{W}_2 \, \text{ReLU}(\mathbf{W}_1 \, \text{GAP}(\mathbf{X})) \right)
    \end{equation}
    \item \textbf{Chú ý không gian (Spatial Squeeze-and-Excitation -- sSE):} Chiếu toàn bộ kênh đặc trưng về 1 bản đồ kích hoạt không gian qua lớp tích chập $1 \times 1$:
    \begin{equation}
    \mathbf{z}_{sSE} = \sigma \left( \text{Conv}_{1 \times 1}(\mathbf{X}) \right)
    \end{equation}
\end{enumerate}
Đặc trưng đầu ra được kết hợp cộng gộp trực tiếp:
\begin{equation}
\mathbf{X}_{scSE} = \mathbf{X} \odot \mathbf{z}_{cSE} + \mathbf{X} \odot \mathbf{z}_{sSE}
\end{equation}
Cơ chế này cho phép mạng vừa tập trung vào ``kênh tần số nào quan trọng'' (qua cSE) vừa tập trung vào ``vùng pixel nào có vết cắt biên'' (qua sSE).

\subsubsection{2.4. Các độ đo đánh giá mô hình đa nhiệm}
\indent\indent Hệ thống sử dụng tập hợp các chỉ số đánh giá toàn diện cho cả hai nhiệm vụ:\vspace{-0.1cm}
\begin{itemize}[label={--}, itemsep=1pt]
    \item \textbf{Pixel Accuracy:} Tỷ lệ phân loại đúng từng pixel trên toàn bộ ảnh tài liệu:
    \begin{equation}
    \text{Pixel Acc} = \frac{TP + TN}{TP + TN + FP + FN}
    \end{equation}
    \item \textbf{Intersection over Union (IoU):} Đo lường mức độ trùng khớp giữa mặt nạ dự đoán $P$ và mặt nạ ground-truth $G$:
    \begin{equation}
    \text{IoU} = \frac{|P \cap G|}{|P \cup G|} = \frac{TP}{TP + FP + FN}
    \end{equation}
    \item \textbf{Macro Dice / Dice Coefficient:}
    \begin{equation}
    \text{Dice} = \frac{2 \times |P \cap G|}{|P| + |G|} = \frac{2 \times TP}{2 \times TP + FP + FN}
    \end{equation}
    \item \textbf{Top-k Class Hit Rate (đặc biệt Top-3 Hit Rate):} Tỷ lệ các mẫu kiểm thử mà trong đó ít nhất một nhãn làm giả thực tế (ground-truth positive) xuất hiện trong top $k$ nhãn có độ tin cậy dự đoán cao nhất:
    \begin{equation}
    \text{Hit@k} = \frac{1}{N} \sum_{i=1}^{N} \mathbb{I} \left( \text{Top-k}(\hat{\mathbf{y}}_i) \cap \mathbf{Y}_i^+ \neq \emptyset \right)
    \end{equation}
    Chỉ số này phản ánh năng lực gợi ý chính xác dấu vết nghi vấn cho chuyên viên giám định.
\end{itemize}
